\scanpage{78}
\mspage{22}

\deleted{\unsure{Comme}} \deleted{Si on a une $W$-équivalence}

Ici encore, la construction est fonctorielle, on trouve le foncteur
\[ C:\mathrm{Cat}_{/S}\longrightarrow\underline{\mathrm{Parf}}\,S \]
qu'il faut comparer au foncteur en sens inverse
\[ r:\mathrm{Cat}_{/S}\hookleftarrow\underline{\mathrm{Parf}}\,S. \]

Tout d'abord, vérifions que $C$, \deleted{\unsure{tout}} comme $r$, est compatible avec les \deleted{flèches} localiseurs.

Soit $f:X\to Y$ dans $\mathrm{Cat}_{/S}$, supposons $f$ est dans $W$, \deleted{alors} \inserted{prouvons $C(\unsure{s|S|f})$ est dans $W_S$.} On a commutativité \unsure{du|des}
\diagram{\begin{tikzcd}[ampersand replacement=\&]
\overset{\deleted{\unsure{C}}}{X}\arrow[r,"f"]\arrow[d,"j_X"']\&Y\arrow[d,"j_Y"]\\
CX\arrow[r,"Cf"']\&CY
\end{tikzcd}}
et $j_X,j_Y$ sont dans $W$, donc $f$ est dans $W$ ssi $Cf$ l'est. \deleted{\unsure{que}} Mais dans le cas $W=W_\infty$, on sait que cela équivaut \unsure{à|:} ce que $Cf$ soit dans $W_\infty$ fibre par fibre. Admettant \textcircled{1} p.~15 pour $W$, il reste à s'assurer \deleted{que} d'une réciproque \unsure{dans|des|sur} $\underline{\mathrm{Parf}}$, i.e.

\textcircled{2} Soit
\diagram{\begin{tikzcd}[ampersand replacement=\&,column sep=tiny,row sep=tiny]
X\arrow[rr,"f"]\arrow[dr,no head]\&\&Y\arrow[dl,no head]\\
\&S\&
\end{tikzcd}}
dans $\underline{\mathrm{Parf}}_{/S}$ (ou seulement dans $\underline{\mathrm{Fib}\,\mathrm{hot}}\,S$). Si $f$ est dans $W$, alors les $f_s:X_s\unsure{\text{---}|\to}Y_s$ aussi.

% archive sheet number: 77
